% !TEX program = xelatex
% =============================================================================
% cn-gongwen · 党政机关公文模板（GB/T 9704—2012《党政机关公文格式》）
%
% 逐条落实标准的版式参数（引用处标注条号）：
%   5.2.1  A4；天头 37mm、订口 28mm，版心 156mm×225mm
%   5.2.2  各要素一般用 3 号仿宋体
%   5.2.3  每面 22 行、每行 28 字，撑满版心（行距由 linespread 推出，见下）
%   7.2    版头：份号/密级/紧急程度、发文机关标志(红色小标宋、上缘距版心上缘
%          35mm)、发文字号(标志下空二行、六角括号)、红色分隔线(字号下 4mm)
%   7.3    主体：标题(2号小标宋、红线下空二行)、主送机关(顶格全角冒号)、
%          正文(段首空二字；层次 一、黑体/（一）楷体/1.（1）仿宋)、附件说明、
%          发文机关署名与成文日期(加盖印章式：日期右空四字、署名居中其上)、
%          附注、附件(另面，"附件1"3号黑体顶格)
%   7.4    版记：粗线0.35mm—抄送(4号仿宋)—细线0.25mm—印发机关和日期—末条粗线
%   7.5    页码：4号半角宋体数字两侧一字线，距版心下缘 7mm，单页右空一字、
%          双页左空一字（twoside）
%
% 编译：bash scripts/compile.sh templates/cn-gongwen/main.tex --preview
% 字体：方正小标宋/仿宋_GB2312/楷体_GB2312 已安装则自动启用，否则回退到
%       ctex 按操作系统选择的黑体/仿宋/楷体近似替代（见 README「字体合规」）。
% =============================================================================
% linespread 推导：版心高 225mm ÷ 22 行 = 10.227mm/行 ≈ 28.98bp；
% 正文 3 号 = 16bp，28.98/16 ≈ 1.81（ctex 的 linespread 为字号倍数；
% 3 号非 zihao 类选项的合法值，故在 \AtBeginDocument 里用 \zihao{3} 设置）。
\documentclass[linespread=1.81, twoside]{ctexart}

% --- 5.2.1 页边与版心 --------------------------------------------------------
% twoside 下 left=订口(内侧)28mm、right=切口(外侧)26mm，双页自动镜像。
% footskip 使页码一字线上缘距版心下缘约 7mm（7.5）。
\usepackage[a4paper, top=37mm, bottom=35mm, left=28mm, right=26mm,
            footskip=10.5mm]{geometry}
\setlength{\parindent}{2\ccwd}   % 正文每自然段左空二字（7.3.3）
\raggedbottom                    % 禁止竖向胶水拉伸，保持 22 行网格与版头定位稳定

% --- 公文红（6.2：红色油墨达到色谱 Y80%、M80%）------------------------------
\usepackage{xcolor}
\definecolor{gwred}{cmyk}{0, 0.8, 0.8, 0}

% --- 字体：正版公文字体自动检测，缺失则回退 ---------------------------------
% 小标宋（发文机关标志、标题）：无免费授权字体，回退为宋体加粗。
\IfFontExistsTF{FZXiaoBiaoSong-B05S}{%
  \setCJKfamilyfont{gwxbs}{FZXiaoBiaoSong-B05S}%
  \newcommand{\xbsong}{\CJKfamily{gwxbs}}%
}{%
  \IfFontExistsTF{方正小标宋简体}{%
    \setCJKfamilyfont{gwxbs}{方正小标宋简体}%
    \newcommand{\xbsong}{\CJKfamily{gwxbs}}%
  }{%
    \newcommand{\xbsong}{\songti\bfseries}%
  }%
}
% 仿宋_GB2312 / 楷体_GB2312（存在则替换 ctex 默认仿宋/楷体）
\IfFontExistsTF{FangSong_GB2312}{%
  \setCJKfamilyfont{gwfs}{FangSong_GB2312}%
  \renewcommand{\fangsong}{\CJKfamily{gwfs}}%
}{}
\IfFontExistsTF{KaiTi_GB2312}{%
  \setCJKfamilyfont{gwkt}{KaiTi_GB2312}%
  \renewcommand{\kaishu}{\CJKfamily{gwkt}}%
}{}
\renewcommand{\familydefault}{\rmdefault}
\AtBeginDocument{\zihao{3}\fangsong}   % 5.2.2 全文默认 3 号仿宋（5.2.3 行距随 linespread）
\xeCJKsetup{CJKecglue={}}              % 公文习惯：汉字与数字/字母间不加空隙（见式样）

% --- 7.3.3 结构层次：一、黑体 /（一）楷体；1. 与（1）仿宋随正文排 -----------
% 标题行不加额外行距，保持 22 行网格。
\ctexset{
  section = {
    name      = {,、},
    number    = \chinese{section},
    format    = \heiti\zihao{3},
    aftername = {},
    indent    = 2\ccwd,
    beforeskip = 0pt,
    afterskip  = 0pt,
    afterindent = true,
  },
  subsection = {
    name      = {（,）},
    number    = \chinese{subsection},
    format    = \kaishu\zihao{3},
    aftername = {},
    indent    = 2\ccwd,
    beforeskip = 0pt,
    afterskip  = 0pt,
    afterindent = true,
  },
}

% --- 7.5 页码：4号半角宋体数字、两侧一字线；单页右空一字、双页左空一字 -------
\usepackage{fancyhdr}
\pagestyle{fancy}
\fancyhf{}
% 一字 = 4号汉字宽 14bp（页脚上下文中 \ccwd 不可靠，用字面量）
\newcommand{\gwpagenum}{{\zihao{4}\songti —\hspace{4.2bp}\thepage\hspace{4.2bp}—}}
\fancyfoot[RO]{\gwpagenum\hspace*{14bp}}
\fancyfoot[LE]{\hspace*{14bp}\gwpagenum}
\renewcommand{\headrulewidth}{0pt}

% --- 可调参数（标准未规定的量）----------------------------------------------
\newlength{\gwlogodrop}\setlength{\gwlogodrop}{35mm}  % 机关标志上缘距版心上缘(7.2.4)
\newlength{\gwredrule}\setlength{\gwredrule}{0.5mm}   % 版头红色分隔线粗细(标准未定)
\newcommand{\gwlogosize}{\zihao{-0}}                  % 机关标志字号(7.2.4"醒目美观")

\usepackage{graphicx}                % 印章图片可用（见署名处注释）
\usepackage[hidelinks]{hyperref}
\hypersetup{pdftitle={示例市人民政府关于加强政务数据共享管理工作的通知}}

\begin{document}

% =============================================================================
% 版头（7.2）
% =============================================================================
% 份号/密级/紧急程度做成零高度叠排，不占版面高度，保证机关标志定位不受
% 增删行数影响。按需保留或注释掉整块。
\noindent\raisebox{0pt}[0pt][0pt]{\parbox[t]{\textwidth}{%
  000001\\                        % 7.2.1 份号：6位3号数字，顶格第一行
  % {\heiti 机密★1年}\\          % 7.2.2 密级和保密期限：3号黑体，第二行
  % {\heiti 特急}\\               % 7.2.3 紧急程度：3号黑体，第三行
  \mbox{}}}\par
% 7.2.4 发文机关标志：上边缘至版心上边缘 35mm，红色小标宋居中
% （\vspace 扣除首行 \topskip、标志行行高贡献与实测残差 5.9mm，
%   按 150dpi 渲染实测校准：标志上缘落在纸面 72mm = 天头 37 + 35）
\vspace*{\dimexpr\gwlogodrop-\topskip-\baselineskip-5.5mm\relax}
{\centering\gwlogosize\xbsong\color{gwred}示例市人民政府文件\par}
% 7.2.5 发文字号：标志下空二行居中；年份全称、六角括号、顺序号不虚位不加"第"
\vspace{2\baselineskip}
{\centering 示政发〔2026〕8号\par}
\nointerlineskip
% 上行文（请示/报告）改用下面一行：发文字号居左空一字、签发人居右空一字，
% "签发人"用3号仿宋、姓名用3号楷体（7.2.5、7.2.6）：
% \noindent\hspace*{1\ccwd}示政发〔2026〕8号\hfill 签发人：{\kaishu 张示例}\hspace*{1\ccwd}
% 7.2.7 版头分隔线：发文字号之下 4mm，与版心等宽，红色
\vspace{4mm}
\noindent{\color{gwred}\rule{\textwidth}{\gwredrule}}\par
\nointerlineskip

% =============================================================================
% 主体（7.3）
% =============================================================================
% 7.3.1 标题：2号小标宋，红线下空二行居中；回行词意完整、梯形或菱形
% （多行标题行距标准未规定，按公文惯例取约 30bp 紧凑排布）
\vspace{2\baselineskip}
{\centering\zihao{2}\xbsong\baselineskip=30bp
 示例市人民政府\\关于加强政务数据共享管理工作的通知\par}
% 7.3.2 主送机关：标题下空一行，居左顶格，全角冒号
\vspace{\baselineskip}
\noindent 各县（区）人民政府，市政府各部门，市直各单位：

为提升政务服务效能，推进政务数据依法有序共享，根据《中华人民共和国数据安全法》
及省有关规定，结合我市实际，现将有关事项通知如下。

\section{总体要求}

坚持统筹协调、按需共享、安全可控的原则，以政务数据目录为基础，以共享平台为
枢纽，推动跨部门、跨层级业务协同，实现群众和企业办事材料明显精简、办理时限
明显压缩。

\section{主要任务}

\subsection{编制政务数据目录}

各部门应当于2026年9月底前编制本部门政务数据目录，明确数据项、更新频率与
共享属性，报市政务服务和数据管理局汇总审核。

\subsection{规范数据共享流程}

数据共享按以下方式办理：1.无条件共享类，通过市共享平台直接获取；2.有条件
共享类，由需求部门提出申请，提供部门应当在5个工作日内答复；不予共享的，
应当说明理由并提供依据。具体情形包括：（1）涉及国家秘密的；（2）法律法规
另有规定的。

\subsection{强化安全管理}

各部门应当落实数据分类分级保护要求，健全访问控制与日志审计机制，签订数据
使用承诺书，确保共享数据不用于约定范围之外的用途。

\section{保障措施}

各县（区）、各部门要将政务数据共享工作纳入年度绩效考核，明确分管负责同志和
联络员。市政务服务和数据管理局要加强督促指导，每季度通报工作进展。执行中的
重大问题，请径向市政府报告。

% 7.3.4 附件说明：正文下空一行、左空二字；名称后不加标点；回行与名称首字对齐
\vspace{\baselineskip}
\noindent\hspace{2\ccwd}附件：1.政务数据共享责任清单\par
\noindent\hspace{2\ccwd}\hphantom{附件：}2.政务数据共享工作联络表\par

% 7.3.5.1 发文机关署名、成文日期（加盖印章式）：
% 成文日期右空四字，发文机关署名在日期之上、以日期为准居中。
% 如需加盖印章（红色、透明底 PNG，端正居中下压署名和日期，7.3.5.1），
% 取消注释并微调宽度与位移：
% \noindent\raisebox{0pt}[0pt][0pt]{\hspace*{\dimexpr\textwidth-70mm\relax}%
%   \raisebox{-32mm}{\includegraphics[width=42mm]{seal.png}}}\par
\vspace{\baselineskip}
\noindent\hfill
\begin{tabular}{@{}c@{}}
  示例市人民政府\\
  2026年7月18日
\end{tabular}\hspace*{4\ccwd}\par
% 7.3.6 附注：成文日期下一行，居左空二字，加圆括号
\noindent\hspace{2\ccwd}（此件公开发布）\par

% =============================================================================
% 附件（7.3.7）：另面编排；"附件+序号"3号黑体顶格第一行，标题居中在第三行
% =============================================================================
\clearpage
\noindent{\heiti 附件1}\par
\vspace{\baselineskip}
{\centering\zihao{2}\xbsong 政务数据共享责任清单\par}
\vspace{\baselineskip}

\begin{center}
  \renewcommand{\arraystretch}{1.6}
  {\zihao{4}
  \begin{tabular}{|c|p{38mm}|p{34mm}|p{28mm}|}
    \hline
    序号 & 数据类别 & 提供部门 & 共享属性 \\
    \hline
    1 & 企业登记基础信息 & 市市场监管局 & 无条件共享 \\
    \hline
    2 & 不动产登记信息 & 市自然资源局 & 有条件共享 \\
    \hline
    3 & 社会保险参保信息 & 市人力资源社会保障局 & 有条件共享 \\
    \hline
    4 & 纳税信用评价信息 & 市税务局 & 有条件共享 \\
    \hline
  \end{tabular}}
\end{center}

各提供部门应当在数据变更之日起3个工作日内完成目录与数据更新，并对数据的
真实性、完整性负责。

% =============================================================================
\clearpage
\noindent{\heiti 附件2}\par
\vspace{\baselineskip}
{\centering\zihao{2}\xbsong 政务数据共享工作联络表\par}
\vspace{\baselineskip}

各县（区）、各部门确定1名分管负责同志和1名联络员，于2026年8月10日前将下列
信息报市政务服务和数据管理局：

——单位名称、分管负责同志姓名及职务；

——联络员姓名、办公电话与政务邮箱；

——本单位政务数据目录编制工作进度安排。

联系人：市政务服务和数据管理局数据共享科；电话：010-8800 0000。

% =============================================================================
% 版记（7.4）：置于公文最后一面版心最下方，末条分隔线与版心下边缘重合
% =============================================================================
\vfill
\begingroup\zihao{4}\setlength{\parindent}{0pt}%
\rule{\textwidth}{0.35mm}\par                    % 首条分隔线（粗，0.35mm）
\vspace{0.25\baselineskip}
% 7.4.2 抄送机关：4号仿宋，左右各空一字，回行与冒号后首字对齐，末尾句号
\hspace{1\ccwd}\begin{minipage}{\dimexpr\textwidth-2\ccwd\relax}
  \hangindent=3\ccwd\hangafter=1\relax
  抄送：市委各部门，市人大常委会办公室，市政协办公室，市监委，市法院，市检察院。\par
\end{minipage}\par
\vspace{0.25\baselineskip}
\rule{\textwidth}{0.25mm}\par                    % 中间分隔线（细，0.25mm）
\vspace{0.2\baselineskip}
% 7.4.3 印发机关左空一字、印发日期右空一字加"印发"
\hspace{1\ccwd}示例市人民政府办公室\hfill 2026年7月18日印发\hspace*{1\ccwd}\par
\vspace{0.2\baselineskip}
\rule{\textwidth}{0.35mm}\par                    % 末条分隔线（粗，与版心下缘重合）
\endgroup

\end{document}
